\section{Method}
\label{sec:method}

\BM keeps \MS's renderer, losses and schedule and changes how the mesh grows and shrinks so that the encoded file meets a target size $T$. The logic follows from the goal. To hit a size we must know what the mesh costs in bytes while it is still being trained (Sec.~\ref{sec:probe}); the cost turns $T$ into a vertex count and a growth plan (Sec.~\ref{sec:growth}); whatever the plan could not foresee is removed by edge collapse during training (Sec.~\ref{sec:trim}). Alg.~\ref{alg:bm} gives the loop and Sec.~\ref{sec:codec} the codec whose bytes are counted.

\subsection{Preliminaries}
\label{sec:prelim}
\MS~\cite{meshsplatting2025} represents a scene by $V$ vertices, each with a position $\mathbf{x}\in\mathbb{R}^3$, a DC color $\mathbf{c}\in\mathbb{R}^3$ and 45 spherical-harmonic (SH) coefficients $\mathbf{h}$ of degrees 1--3, and $F$ triangles over them. Training runs 30k iterations. Every 500 iterations until 15k it \emph{densifies}: it samples triangles by an importance score and replaces each by four, inserting the three edge midpoints as new vertices; it also prunes triangles that render with low opacity. A vertex cap stops densification once $V$ reaches the cap. At 11k a restricted Delaunay step stitches the soup into a connected mesh, an opacity floor then rises until every triangle is opaque, and a final cleaning removes triangles that cover no pixel. We keep all of this. Let $R$ be the size of the encoded file in bytes. We write $D_t$ for the \emph{error of triangle} $t$: the per-pixel $L_1$ residual summed over the pixels whose front-most triangle is $t$ (the rasterizer returns this id), averaged over the views rendered since the last decision and normalized per megapixel.

\subsection{Measuring the file during training}
\label{sec:probe}
\noindent\textbf{Byte probe.} Every 1000 iterations once \MS starts pruning (4k), a copy of the mesh is cleaned exactly as the final mesh will be and passed through the real encoder of Sec.~\ref{sec:codec}. The file size is affine in the mesh,
\begin{equation}
R = b_v V + b_f F + c ,
\label{eq:bytes}
\end{equation}
where $c$ covers the dictionaries and the header. The probe returns $R$, the bytes per vertex $b_v$ and per face $b_f$, $c$, and the fraction of vertices that survive cleaning; these replace the previous values. Measuring matters because the costs are not constants: $b_v$ ranges from 8 to 13 bytes across scenes and operating points, and with fixed costs small files land 18\% too large (Sec.~\ref{sec:e7}).

\noindent\textbf{Vertex count.} With $f$ faces per vertex, Eq.~\eqref{eq:bytes} gives the vertex count that the target buys:
\begin{equation}
V^\star = (\alpha T - c)\,/\,(b_v + f\, b_f).
\label{eq:vstar}
\end{equation}
During growth the mesh is an unstitched soup whose face-to-vertex ratio says little about the final mesh, so we fix $f = 2.5$; the factor $\alpha$ is discussed in Sec.~\ref{sec:trim}.

\subsection{Growing toward the vertex count}
\label{sec:growth}
\noindent\textbf{How many.} With $m$ growth steps left, the plan raises $V$ geometrically, $V \rightarrow V\,(V^\star/V)^{1/m}$, which asks for
\begin{equation}
k = V\big((V^\star/V)^{1/m}-1\big)\,/\,(3\sigma)
\label{eq:k}
\end{equation}
splits now: each split adds three vertices, of which a fraction $\sigma$ survives \MS's pruning at the next step ($\sigma$ is measured from the previous cohort). The plan is recomputed at every step with the latest probe, so an error in one step is spread over the steps that remain.

\noindent\textbf{Which.} \MS samples the triangles to split by an importance score that favours visible triangles. We split the $k$ triangles with the largest error $D_t$: a triangle is refined because it is rendered wrongly. Written as a threshold, triangle $t$ is split iff
\begin{equation}
D_t \,/\, \Delta R_{\mathrm{split}} > \lambda ,\qquad \Delta R_{\mathrm{split}} = 3(b_v + b_f),
\label{eq:split}
\end{equation}
with $\lambda$ set to the $k$-th largest ratio. A split costs the same bytes for every triangle, so Eq.~\eqref{eq:split} is a ranking by error and $\lambda$ is a result of the plan, not a parameter; we keep the form because the same quantity, error per byte, also decides collapses (Sec.~\ref{sec:trim}), where the bytes do differ. Because the errors are re-measured between steps, the number of triangles that actually passes differs from $k$; we divide $k$ by the measured yield $y$ (passed over asked) of the last steps before taking the rank. Running Eq.~\eqref{eq:split} with a fixed $\lambda$ instead gives training without a size target, which serves as the reference in Sec.~\ref{sec:e3}; one such $\lambda$ lands at files $2.5\times$ apart across scenes.

\subsection{Trimming by collapse}
\label{sec:trim}
\noindent\textbf{What growth cannot see.} After the mesh is stitched at 11k, the set of visible faces keeps changing as the opacity floor rises, and the stitched mesh has its own faces per vertex (2.0 to 3.5 on our scenes), none of which exists before 11k. The size reached by growth alone is therefore off by a scene-dependent amount, and the remainder is corrected on the stitched mesh.

\noindent\textbf{Collapse rule.} Collapsing edge $(u,v)$ merges $u$ and $v$ into one vertex and deletes the faces on the edge. It frees $\Delta R_{uv} = b_v\,[\,u,v \text{ survive cleaning}\,] + b_f\, n_{uv}$ bytes, with $n_{uv}$ the deleted faces that would survive cleaning, and it risks the error $D_{uv}$ of the faces around $u$ and $v$. Edges are collapsed in increasing order of $D_{uv}/\Delta R_{uv}$, as an independent set (no two collapses share a vertex) with a fold-over check. The merged vertex is placed at the midpoint and takes the average of the two attribute vectors and of their optimizer moments, so training continues without a restart and the survivors keep being optimized.

\noindent\textbf{When and how much.} From 12.5k to 19.5k, every 500 iterations, the probe measures $R$ and a least-squares fit of the probe history extrapolates it to the end of the opacity ramp. If the estimate exceeds $T$, edges are collapsed until the excess is removed, at most 15\% of the vertices per round. If it does not, nothing is collapsed: trimming removes bytes and cannot add them.

\noindent\textbf{Growth aim.} Since trimming is one-sided, the plan aims above the target by the factor $\alpha$ of Eq.~\eqref{eq:vstar}. $\alpha$ stands for one number per scene that growth cannot observe, the faces per vertex of the final mesh: over 18 runs without trimming, the aim that would have landed exactly on the target correlates with it at $r=-0.97$, and no quantity available during growth predicts it on both captured scenes and objects (supplement). We therefore use one value per capture class, 1.25 for captured scenes and 1.1 for synthetic objects, chosen on Mip-NeRF~360 and NeRF-Synthetic, and let trimming absorb the per-scene remainder. The other settings are fixed for all runs: $f=2.5$ during growth, a probe every 1000 iterations, trimming rounds as above; the initial byte costs are overwritten by the first probe.

\begin{algorithm}[t]
\caption{\BM training for a target size $T$}
\label{alg:bm}
\small
\begin{algorithmic}[1]
\Require images, target $T$, aim $\alpha$
\State initialize \MS; byte costs $(b_v, b_f, c)$ from defaults
\For{iteration $i = 1 \dots 30\text{k}$}
  \State \MS optimization step; accumulate $D_t$ \Comment{per-triangle error}
  \If{growth step ($i \le 15\text{k}$)}
    \State probe every 1k: encode a cleaned copy $\rightarrow$ $b_v, b_f, c$
    \State $V^\star \gets (\alpha T - c)/(b_v + f b_f)$ \Comment{Eq.~\ref{eq:vstar}}
    \State $k \gets$ geometric plan toward $V^\star$, survival $\sigma$ \Comment{Eq.~\ref{eq:k}}
    \State split the $k/y$ triangles with the largest $D_t$ \Comment{Eq.~\ref{eq:split}}
  \EndIf
  \If{trimming step ($12.5\text{k} \le i \le 19.5\text{k}$)}
    \State $\hat R \gets$ probe, extrapolated to end of ramp
    \If{$\hat R > T$}
      \State collapse lowest $D_{uv}/\Delta R_{uv}$ until $\hat R \le T$
    \EndIf
  \EndIf
\EndFor
\State final cleaning; encode (Sec.~\ref{sec:codec}) $\rightarrow$ one file
\end{algorithmic}
\end{algorithm}

\subsection{Codec}
\label{sec:codec}
The trained mesh stores 52 floats per vertex and three 32-bit indices per face. The codec makes three lossy reductions, then codes the result losslessly in an order where each stream is predicted from what the decoder already has.

\noindent\textbf{Lossy reductions.} (i)~Positions are quantized on a per-axis uniform grid of $b=\lceil\log_2(\text{extent}/(0.1\,\bar e))\rceil$ bits, $\bar e$ the median edge length (14--18 bits on our scenes): the step is a tenth of a typical edge; one bit finer costs bytes for no PSNR, one bit coarser loses 0.06--0.2\dB. (ii)~The 45 SH coefficients are 87\% of the floats, and neighbouring vertices on a surface share their view-dependent behaviour, but the coefficients are too noisy to predict from one another. We therefore replace them by a two-level per-scene dictionary, $\mathbf{h}_v \approx \mathbf{B}[i_v] + \mathbf{B}_r[j_v]$, with $|\mathbf{B}|=4096$ and $|\mathbf{B}_r|=256$ entries fitted by $k$-means and refined for 400 steps against the rendering loss with the assignments fixed. Each vertex keeps two indices. (iii)~The DC color is quantized to 7 bits per channel, and the opacity, uniform after training, is one scalar.

\noindent\textbf{Vertex order.} Faces are reordered into a cache-friendly strip order~\cite{meshoptimizer}, and vertex ids are assigned in order of first reference. The decoder can then infer every id: a vertex that has not appeared yet is always the next id. No permutation is transmitted, and every later stream can use the already-decoded neighbours of a vertex.

\noindent\textbf{Connectivity.} The coder keeps the 64 most recently decoded directed edges. A face that shares an edge with one of them is coded as \emph{hit}: the cache slot, the edge direction and the third vertex; otherwise as \emph{miss}: three vertices. A vertex is coded as 0 when it is new and as the difference to the last coded id otherwise. All symbols are coded with an adaptive binary arithmetic coder whose contexts are the symbol type and the bits already coded; no tables are sent. Edgebreaker-style coders traverse a manifold; our meshes are stitched soups with 45\% boundary edges and thousands of components, which forces Edgebreaker to split vertices, so a local cache is the better model here. Connectivity costs about 15 bits per face and is the largest stream (39--43\% of the file).

\noindent\textbf{Positions.} Each quantized position is predicted by its earliest already-decoded mesh neighbour, rebuilt by the decoder from the faces, and the residual is arithmetic coded per axis with the magnitude of the previous residual as context: about 16 bits per vertex against 45 raw.

\noindent\textbf{Appearance indices.} Index frequencies are nearly uniform, so an order-0 entropy coder gains nothing; the redundancy is spatial. For vertex $v$ the coder lists the distinct indices of its already-decoded neighbours and codes one flag (is $i_v$ among them) and, if so, its position in the list; otherwise the index itself. The DC colors are byte-shuffled and compressed with zstd.

\noindent\textbf{Decoding.} The decoder reads the connectivity, rebuilds the neighbour graph from it, then decodes positions and indices and expands $\mathbf{h}_v$ from the dictionaries. It needs nothing but the file, runs in 2--6~s on one CPU core, and outputs a standard indexed mesh that an ordinary vertex shader draws. The encoder verifies a byte-exact round trip before writing; all loss comes from the three reductions above.
